\begin{theorem}[Forward simulation is a congruence for the composition operators]
  \label{thm:forwardsim-congruence}
  \lean{PLTS.ForwardSimulation.parallel_right}\lean{PLTS.ForwardSimulation.parallel_left}\lean{PLTS.ForwardSimulation.synchronisedProduct}\lean{PLTS.ForwardSimulation.abstract}\lean{PLTS.ForwardSimulation.relabel}\lean{PLTS.ForwardSimulation.trans}\leanok
  \uses{def:forwardsim, def:process-algebra, def:relabel, def:syncproduct}
  Forward simulation between labelled transition systems survives each operator a composition is built from --- parallel composition in either position, the full-synchronisation product of a finite family, hiding a label set, and restriction along the left embedding --- and two forward simulations compose: \leandecl{PLTS.ForwardSimulation.parallel_right}, \leandecl{PLTS.ForwardSimulation.parallel_left}, \leandecl{PLTS.ForwardSimulation.synchronisedProduct}, \leandecl{PLTS.ForwardSimulation.abstract}, \leandecl{PLTS.ForwardSimulation.relabel} and \leandecl{PLTS.ForwardSimulation.trans}.
  Replacing one component of a composition by a system that simulates it is an application of these, and every substitution inside a gather instance and inside a graded-agreement round is proved that way.
\end{theorem}

\begin{proof}\leanok
  \uses{def:forwardsim, def:process-algebra, def:syncproduct, def:weakstep}
  See the Lean proofs of \leandecl{PLTS.ForwardSimulation.parallel_right}, \leandecl{PLTS.ForwardSimulation.synchronisedProduct}, \leandecl{PLTS.ForwardSimulation.abstract}, \leandecl{PLTS.ForwardSimulation.relabel} and \leandecl{PLTS.ForwardSimulation.trans}.
\end{proof}
